\documentclass[12pt]{article}

\usepackage[margin=1in]{geometry}
\usepackage{amsmath, amssymb, amsthm}
\usepackage[T1]{fontenc}
\usepackage{lmodern}
\usepackage{microtype}
\usepackage{setspace}
\usepackage{booktabs}
\usepackage{array}
\usepackage{caption}
\usepackage[colorlinks=true, allcolors=blue]{hyperref}

\setlength{\parindent}{0pt}
\setlength{\parskip}{0.6em}
\onehalfspacing

\captionsetup{
    font=small,
    labelfont=bf,
    justification=centering
}

\newtheorem{theorem}{Theorem}[section]
\newtheorem{lemma}[theorem]{Lemma}
\newtheorem{proposition}[theorem]{Proposition}
\newtheorem{corollary}[theorem]{Corollary}

\theoremstyle{definition}
\newtheorem{definition}[theorem]{Definition}

\theoremstyle{remark}
\newtheorem{remark}[theorem]{Remark}

\title{\textbf{Scott's Electrical Variation Response Equation}}

\author{
Jaxson Scott Levoy Freeman\\
\small Independent student researcher\\
\small Corresponding author: \texttt{jaxsonstormy@gmail.com}
}

\date{November 2026}

\begin{document}

\maketitle

\begin{abstract}

This paper introduces and investigates a parameterized
mathematical functional designed to represent an
accumulated electrical response to historical voltage
variations.

The proposed Electrical Variation Response Equation
(EVRE) is defined by

\[
\Omega_s(t)=
\sqrt{
\Omega_0^2+
\frac{\beta}{T}
\int_0^t
e^{-\lambda(t-\tau)}
\left(\frac{dV(\tau)}{d\tau}\right)^2
\,d\tau
},
\]

where \(\Omega_0\) represents an initial response,
\(V(t)\) denotes a time-dependent voltage signal,
\(\beta\) is a positive scaling parameter,
\(T\) is a reference time interval, and
\(\lambda\) is a positive memory-decay parameter.

The proposed functional combines squared voltage
variation, exponential temporal weighting, and
a nonlinear square-root transformation.

We investigate its dimensional consistency,
nonnegativity, boundedness under specified input
conditions, continuity, and dependence on its
parameters.

An equivalent first-order differential equation
is derived, allowing the historical integral to
be represented through a local evolution equation.
We also examine the behavior under constant voltage,
scaled voltage inputs, and bounded voltage variation.

A reconstruction expression is obtained for the
magnitude of the instantaneous voltage derivative
when the response and model parameters are known.

The proposed equation is a mathematical construction
rather than an established electrical law. Its
physical interpretation, practical usefulness, and
empirical significance require further investigation.

\end{abstract}

\noindent\textbf{Keywords:}
Electrical response; voltage variation;
exponential memory; integral functionals;
nonlinear mathematical modeling; electrical systems;
parameter sensitivity; differential equations.

\section{Introduction}

Electrical systems are commonly investigated
through relationships involving voltage, current,
resistance, capacitance, and inductance.

Many electrical quantities are described using
instantaneous relationships between circuit
variables. However, mathematical models can also
be constructed to investigate quantities that
depend on the previous behavior of a signal.

A time-dependent voltage signal may experience
continuous variation during the operation of
an electrical system.

The magnitude and frequency of these variations
can be investigated mathematically through
derivatives and integral expressions.

This paper introduces a parameterized functional
that accumulates squared voltage variation over
a historical time interval while assigning
exponentially decreasing weight to older changes.

The resulting expression introduces a
memory-dependent mathematical response quantity
denoted by \(\Omega_s(t)\).

The purpose of this investigation is to establish
the mathematical properties of the proposed
functional and determine whether it provides
a useful framework for describing accumulated
electrical variation.

The principal questions considered are:

\begin{enumerate}
    \item Is the proposed functional mathematically
    well-defined and dimensionally consistent?

    \item Is its response always nonnegative?

    \item Under what conditions is the response bounded?

    \item How does the memory parameter influence
    the contribution of historical voltage changes?

    \item Can the integral representation be
    converted into a differential equation?

    \item Is information about the instantaneous
    voltage variation recoverable from the response?
\end{enumerate}

The investigation begins with the definition
of the mathematical quantities and proceeds
toward the construction of the proposed equation.

\section{Mathematical Foundations}

\subsection{The voltage function}

Let

\[
V:[0,\infty)\rightarrow\mathbb{R}
\]

represent a time-dependent voltage signal.

The voltage is assumed to be locally absolutely
continuous, ensuring that its derivative exists
almost everywhere and that the corresponding
integral expressions can be investigated.

Define the instantaneous voltage variation by

\begin{equation}
v'(t)=\frac{dV(t)}{dt}.
\label{eq:voltage-derivative}
\end{equation}

The derivative measures the rate at which
the voltage changes with respect to time.

Its units are

\[
[v'(t)]=\frac{\mathrm{V}}{\mathrm{s}}.
\]

The square of this quantity is used in the
construction of the proposed functional.

\subsection{The historical weighting function}

Introduce an exponential weighting function:

\begin{equation}
K(t,\tau)=e^{-\lambda(t-\tau)},
\label{eq:kernel}
\end{equation}

where

\[
\lambda>0
\]

is the memory-decay parameter.

For

\[
0\leq\tau\leq t,
\]

the weighting function satisfies

\[
0<K(t,\tau)\leq1.
\]

At the present time, where \(\tau=t\),

\[
K(t,t)=1.
\]

For earlier times, the weighting function
decreases as the difference \(t-\tau\) increases.

Thus, recent voltage variations receive greater
weight than older voltage variations.

\subsection{The accumulated variation functional}

Define the accumulated historical variation
by

\begin{equation}
Q(t)=
\int_0^t
e^{-\lambda(t-\tau)}
\left(\frac{dV(\tau)}{d\tau}\right)^2
\,d\tau.
\label{eq:variation-functional}
\end{equation}

This quantity represents an exponentially
weighted accumulation of squared voltage
variation.

Since the integrand is nonnegative,

\[
Q(t)\geq0.
\]

The functional \(Q(t)\) provides the mathematical
foundation for the proposed response equation.

\section{Definition of the Electrical Variation Response Equation}

\subsection{Construction}

The accumulated variation functional provides
a nonnegative quantity that depends on the
historical behavior of a voltage signal.

To construct a response quantity with units
of electrical resistance, introduce the following
parameters:

\begin{itemize}
    \item \(\Omega_0\): Initial response value.
    \item \(\beta\): Positive scaling coefficient.
    \item \(T\): Positive reference time interval.
    \item \(\lambda\): Positive memory-decay parameter.
\end{itemize}

The parameters satisfy

\[
\Omega_0\geq0,\qquad
\beta>0,\qquad
T>0,\qquad
\lambda>0.
\]

\subsection{Formal definition}

\begin{definition}[Electrical Variation Response Equation]

For an admissible voltage signal \(V(t)\),
define the Electrical Variation Response Equation
by

\begin{equation}
\Omega_s(t)=
\sqrt{
\Omega_0^2+
\frac{\beta}{T}
\int_0^t
e^{-\lambda(t-\tau)}
\left(\frac{dV(\tau)}{d\tau}\right)^2
\,d\tau
}.
\label{eq:evre}
\end{equation}

\end{definition}

The quantity \(\Omega_s(t)\) is defined as a
memory-dependent electrical response measure.

It is not automatically equivalent to the
physical resistance of a component.

Its interpretation as a resistance-related
quantity would require additional physical
justification.

\subsection{Dimensional consistency}

The proposed expression must be dimensionally
consistent.

The voltage derivative has dimensions

\[
\left[\frac{dV}{dt}\right]
=
\frac{\mathrm{V}}{\mathrm{s}}.
\]

Therefore,

\[
\left[
\left(\frac{dV}{dt}\right)^2dt
\right]
=
\frac{\mathrm{V}^2}{\mathrm{s}}.
\]

The exponential function is dimensionless
provided that

\[
[\lambda]=\mathrm{s}^{-1}.
\]

Since the quantity inside the square root
must have units of resistance squared,
the scaling parameter must satisfy

\begin{equation}
[\beta]
=
\frac{\Omega^2\mathrm{s}^2}{\mathrm{V}^2}.
\label{eq:beta-units}
\end{equation}

The reference time has units

\[
[T]=\mathrm{s}.
\]

Consequently,

\[
\left[
\frac{\beta}{T}Q(t)
\right]
=
\Omega^2.
\]

Both terms inside the square root therefore
have identical dimensions.

The proposed equation is dimensionally
consistent under these parameter definitions.

\subsection{Equivalent representation}

Using Equation~\eqref{eq:variation-functional},
the response equation can be written as

\begin{equation}
\Omega_s(t)=
\sqrt{
\Omega_0^2+\frac{\beta}{T}Q(t)
}.
\label{eq:evre-q}
\end{equation}

This representation will be useful when
investigating the mathematical properties
of the functional.

\section{Basic Mathematical Properties}

\subsection{Nonnegativity}

\begin{theorem}

For every admissible voltage signal and
all permitted parameter values,

\begin{equation}
\Omega_s(t)\geq\Omega_0\geq0.
\label{eq:nonnegative}
\end{equation}

\end{theorem}

\begin{proof}

From Equation~\eqref{eq:variation-functional},

\[
Q(t)=
\int_0^t
e^{-\lambda(t-\tau)}
\left(V'(\tau)\right)^2d\tau.
\]

The exponential weighting function is positive,
and the squared voltage derivative is
nonnegative.

Therefore,

\[
Q(t)\geq0.
\]

Since

\[
\beta>0,\qquad T>0,
\]

we obtain

\[
\frac{\beta}{T}Q(t)\geq0.
\]

It follows that

\[
\Omega_0^2+\frac{\beta}{T}Q(t)
\geq\Omega_0^2.
\]

Taking nonnegative square roots proves

\[
\Omega_s(t)\geq\Omega_0.
\]

\end{proof}

\subsection{Initial condition}

\begin{proposition}

The response equation satisfies

\begin{equation}
\Omega_s(0)=\Omega_0.
\end{equation}

\end{proposition}

\begin{proof}

At \(t=0\), the integral has identical
upper and lower limits.

Therefore,

\[
Q(0)=0.
\]

Substituting into Equation~\eqref{eq:evre-q}
gives

\[
\Omega_s(0)=\sqrt{\Omega_0^2}.
\]

Since \(\Omega_0\geq0\),

\[
\Omega_s(0)=\Omega_0.
\]

\end{proof}

\subsection{Equality condition}

\begin{proposition}

For a fixed time \(t>0\),

\[
\Omega_s(t)=\Omega_0
\]

if and only if

\[
V'(\tau)=0
\]

for almost every \(\tau\in[0,t]\).

\end{proposition}

\begin{proof}

The equality

\[
\Omega_s(t)=\Omega_0
\]

is equivalent to

\[
Q(t)=0.
\]

Because the weighting function is strictly
positive and the squared derivative is
nonnegative, the integral can equal zero
only when

\[
V'(\tau)=0
\]

almost everywhere on the integration interval.

The converse follows immediately by substitution.

\end{proof}

\subsection{Continuity}

\begin{theorem}

If \(V'\) is locally square-integrable, then
\(\Omega_s(t)\) is continuous on every finite
time interval.

\end{theorem}

\begin{proof}

Let

\[
f(\tau)=\left(V'(\tau)\right)^2.
\]

By assumption, \(f\) is locally integrable.

Rewrite the accumulated functional as

\[
Q(t)=
e^{-\lambda t}
\int_0^t e^{\lambda\tau}f(\tau)d\tau.
\]

The integrand is locally integrable.

Its indefinite integral is therefore
absolutely continuous on every finite interval.

Multiplication by the continuous function
\(e^{-\lambda t}\) preserves continuity.

Thus, \(Q(t)\) is continuous.

Finally,

\[
\Omega_s(t)=
\sqrt{\Omega_0^2+\frac{\beta}{T}Q(t)}
\]

is a continuous composition of continuous
functions.

Therefore, the response is continuous.

\end{proof}

\section{Differential Representation}

\subsection{Derivation of the differential equation}

The integral representation contains
information about the entire voltage history.

However, it can also be expressed through
a first-order differential equation.

Recall that

\[
Q(t)=
e^{-\lambda t}
\int_0^t e^{\lambda\tau}
\left(V'(\tau)\right)^2d\tau.
\]

Differentiate using the product rule
and the fundamental theorem of calculus.

We obtain

\[
Q'(t)=
-\lambda e^{-\lambda t}
\int_0^t e^{\lambda\tau}
\left(V'(\tau)\right)^2d\tau
+
\left(V'(t)\right)^2.
\]

Recognizing the original integral gives

\begin{equation}
Q'(t)=
\left(V'(t)\right)^2-\lambda Q(t).
\label{eq:q-ode}
\end{equation}

The initial condition is

\[
Q(0)=0.
\]

This differential equation is equivalent
to the original historical integral under
the stated assumptions.

\subsection{Response evolution}

From

\[
\Omega_s^2(t)=
\Omega_0^2+\frac{\beta}{T}Q(t),
\]

differentiate both sides:

\[
2\Omega_s(t)\Omega_s'(t)
=
\frac{\beta}{T}Q'(t).
\]

Substituting Equation~\eqref{eq:q-ode} gives

\[
2\Omega_s\Omega_s'
=
\frac{\beta}{T}
\left[
(V')^2-\lambda Q
\right].
\]

Since

\[
Q=
\frac{T}{\beta}
\left(\Omega_s^2-\Omega_0^2\right),
\]

we obtain

\begin{equation}
\Omega_s'(t)=
\frac{
\beta(V'(t))^2
-
\lambda T\left(\Omega_s^2(t)-\Omega_0^2\right)
}{
2T\Omega_s(t)
}.
\label{eq:response-ode}
\end{equation}

This equation holds wherever the response
is differentiable and strictly positive.

The differential representation shows that
the response is influenced by two competing
contributions.

The first is the instantaneous voltage
variation.

The second is the decay of accumulated
historical variation.

\subsection{Interpretation of the differential representation}

The term

\[
\beta(V'(t))^2
\]

represents the instantaneous contribution
from voltage variation.

The term

\[
\lambda T
\left(\Omega_s^2(t)-\Omega_0^2\right)
\]

represents the decay of the accumulated
response contribution.

Consequently, the response does not necessarily
increase at every moment.

It can decrease when historical variation
decays faster than new variation is introduced.

\section{Boundedness and Input Conditions}

\subsection{Bounded voltage derivative}

Consider an admissible voltage signal satisfying

\begin{equation}
|V'(t)|\leq M
\end{equation}

for some finite constant \(M\geq0\).

Then

\[
(V'(t))^2\leq M^2.
\]

Substitution into the accumulated functional gives

\[
Q(t)\leq
M^2\int_0^t e^{-\lambda(t-\tau)}d\tau.
\]

Evaluating the integral,

\[
Q(t)\leq
\frac{M^2}{\lambda}
\left(1-e^{-\lambda t}\right).
\]

Therefore, the response satisfies the following
result.

\begin{theorem}

If the voltage derivative is bounded by \(M\),
then

\begin{equation}
\Omega_0
\leq\Omega_s(t)
\leq
\sqrt{
\Omega_0^2+
\frac{\beta M^2}{T\lambda}
\left(1-e^{-\lambda t}\right)
}.
\label{eq:bounded-response}
\end{equation}

\end{theorem}

\begin{proof}

The previous inequality establishes that

\[
0\leq Q(t)\leq
\frac{M^2}{\lambda}
\left(1-e^{-\lambda t}\right).
\]

Multiplying by \(\beta/T\) and adding
\(\Omega_0^2\) gives

\[
\Omega_0^2
\leq
\Omega_s^2(t)
\leq
\Omega_0^2+
\frac{\beta M^2}{T\lambda}
\left(1-e^{-\lambda t}\right).
\]

Taking square roots proves the result.

\end{proof}

\subsection{Long-time boundedness}

Taking the limit as \(t\to\infty\) in
Equation~\eqref{eq:bounded-response} gives

\begin{equation}
\limsup_{t\to\infty}\Omega_s(t)
\leq
\sqrt{
\Omega_0^2+
\frac{\beta M^2}{T\lambda}
}.
\label{eq:long-bound}
\end{equation}

Thus, bounded voltage variation produces
a uniformly bounded response.

\subsection{Constant voltage}

Suppose that the voltage remains constant:

\[
V(t)=V_0.
\]

Then

\[
V'(t)=0.
\]

Consequently,

\[
Q(t)=0
\]

and

\begin{equation}
\Omega_s(t)=\Omega_0.
\end{equation}

A constant voltage therefore produces
no additional accumulated response.

\section{Parameter Sensitivity}

The response depends on three principal
model parameters: \(\beta\), \(T\), and
\(\lambda\).

Their effects are examined independently.

\subsection{Sensitivity with respect to \(\beta\)}

Differentiating Equation~\eqref{eq:evre-q}
with respect to \(\beta\) gives

\begin{equation}
\frac{\partial\Omega_s}{\partial\beta}
=
\frac{Q(t)}
{2T\Omega_s(t)}.
\label{eq:beta-sensitivity}
\end{equation}

For positive accumulated variation,

\[
Q(t)>0.
\]

Therefore,

\[
\frac{\partial\Omega_s}{\partial\beta}>0.
\]

Increasing \(\beta\) increases the modeled
response whenever voltage variation has occurred.

\subsection{Sensitivity with respect to \(T\)}

Differentiating with respect to \(T\) gives

\begin{equation}
\frac{\partial\Omega_s}{\partial T}
=
-\frac{\beta Q(t)}
{2T^2\Omega_s(t)}.
\label{eq:T-sensitivity}
\end{equation}

For positive accumulated variation,

\[
\frac{\partial\Omega_s}{\partial T}<0.
\]

Thus, increasing the reference interval
decreases the modeled response when the
other parameters and input history remain fixed.

\subsection{Sensitivity with respect to \(\lambda\)}

Differentiate the historical functional:

\[
Q(t)=
\int_0^t
e^{-\lambda(t-\tau)}
(V'(\tau))^2d\tau.
\]

Using differentiation under the integral sign,

\[
\frac{\partial Q}{\partial\lambda}
=
-\int_0^t
(t-\tau)e^{-\lambda(t-\tau)}
(V'(\tau))^2d\tau.
\]

Therefore,

\begin{equation}
\frac{\partial\Omega_s}{\partial\lambda}
=
-\frac{\beta}{2T\Omega_s}
\int_0^t
(t-\tau)e^{-\lambda(t-\tau)}
(V'(\tau))^2d\tau.
\label{eq:lambda-sensitivity}
\end{equation}

The integrand is nonnegative.

Consequently,

\[
\frac{\partial\Omega_s}{\partial\lambda}
\leq0.
\]

Increasing the memory-decay parameter
decreases the contribution of historical
voltage variation.

\subsection{Parameter interpretation}

The three parameters influence different
aspects of the proposed model.

\begin{itemize}

    \item \(\beta\) controls the magnitude
    of the accumulated contribution.

    \item \(T\) establishes a reference
    time scale for the response.

    \item \(\lambda\) controls the rate at
    which historical contributions lose weight.

\end{itemize}

These parameters provide independent
mechanisms for adjusting the mathematical
behavior of the functional.

\section{Response to Particular Voltage Signals}

\subsection{Constant voltage slope}

Consider a voltage signal of the form

\[
V(t)=V_0+at,
\]

where \(a\) is a constant voltage slope.

Then

\[
V'(t)=a.
\]

Substituting into Equation~\eqref{eq:variation-functional},

\[
Q(t)=
a^2\int_0^t e^{-\lambda(t-\tau)}d\tau.
\]

Therefore,

\[
Q(t)=
\frac{a^2}{\lambda}
\left(1-e^{-\lambda t}\right).
\]

The response becomes

\begin{equation}
\Omega_s(t)=
\sqrt{
\Omega_0^2+
\frac{\beta a^2}{T\lambda}
\left(1-e^{-\lambda t}\right)
}.
\label{eq:linear-voltage}
\end{equation}

As time increases, the response approaches

\begin{equation}
\lim_{t\to\infty}\Omega_s(t)
=
\sqrt{
\Omega_0^2+
\frac{\beta a^2}{T\lambda}
}.
\end{equation}

\subsection{Sinusoidal voltage}

Consider a sinusoidal voltage signal:

\[
V(t)=A\sin(\omega t).
\]

Its derivative is

\[
V'(t)=A\omega\cos(\omega t).
\]

Therefore,

\[
(V'(t))^2=
A^2\omega^2\cos^2(\omega t).
\]

Using the identity

\[
\cos^2(\omega t)
=
\frac{1+\cos(2\omega t)}{2},
\]

the accumulated functional becomes

\[
Q(t)=
\frac{A^2\omega^2}{2}
\int_0^t
e^{-\lambda(t-\tau)}
\left[1+\cos(2\omega\tau)\right]d\tau.
\]

Evaluating the integral gives

\begin{align}
Q(t)=\frac{A^2\omega^2}{2}
\Bigg[
&\frac{1-e^{-\lambda t}}{\lambda}
\\
&+
\frac{
\lambda\cos(2\omega t)
+2\omega\sin(2\omega t)
-\lambda e^{-\lambda t}
}{
\lambda^2+4\omega^2
}
\Bigg].
\label{eq:sinusoidal-response}
\end{align}

Substitution into the original response
equation provides the corresponding
sinusoidal response.

This demonstrates that the proposed functional
can be evaluated for different classes of
time-dependent voltage signals.

\section{Scaling Properties}

\subsection{Voltage amplitude scaling}

Consider a scaled voltage signal

\[
\widetilde{V}(t)=aV(t),
\]

where \(a\) is a real constant.

Then

\[
\widetilde{V}'(t)=aV'(t).
\]

The accumulated functional becomes

\[
\widetilde{Q}(t)=a^2Q(t).
\]

Consequently,

\begin{equation}
\widetilde{\Omega}_s(t)=
\sqrt{
\Omega_0^2+
a^2\left(\Omega_s^2(t)-\Omega_0^2\right)
}.
\label{eq:scaling}
\end{equation}

Thus, the accumulated contribution scales
quadratically with voltage amplitude before
the square-root transformation.

\subsection{Voltage offset invariance}

Consider adding a constant voltage offset:

\[
\widetilde{V}(t)=V(t)+C.
\]

Differentiation gives

\[
\widetilde{V}'(t)=V'(t).
\]

Therefore,

\begin{equation}
\widetilde{\Omega}_s(t)=\Omega_s(t).
\end{equation}

The proposed functional is invariant under
constant voltage offsets.

Only changes in voltage contribute to
the accumulated response.

\section{Reconstruction of Voltage Variation}

\subsection{Derivation}

The differential representation provides
a method for recovering the magnitude of
the instantaneous voltage derivative.

Starting from Equation~\eqref{eq:response-ode},

\[
2T\Omega_s\Omega_s'
=
\beta(V')^2
-
\lambda T(\Omega_s^2-\Omega_0^2).
\]

Rearranging,

\[
\beta(V')^2
=
2T\Omega_s\Omega_s'
+
\lambda T(\Omega_s^2-\Omega_0^2).
\]

Therefore,

\begin{equation}
|V'(t)|=
\sqrt{
\frac{
2T\Omega_s(t)\Omega_s'(t)
+
\lambda T\left(\Omega_s^2(t)-\Omega_0^2\right)
}{\beta}
}.
\label{eq:reconstruction}
\end{equation}

This expression is valid wherever the
required derivatives exist and the expression
under the square root is nonnegative.

\subsection{Limitations of reconstruction}

The reconstruction expression recovers
the magnitude of the voltage derivative.

It does not independently recover its sign.

Consequently, the response alone cannot
distinguish between equal-magnitude
positive and negative voltage variations.

This is an important limitation of the
proposed functional.

\section{Mathematical Interpretation and Potential Applications}

\subsection{Mathematical interpretation}

The Electrical Variation Response Equation
combines three mathematical operations.

First, the voltage derivative measures
instantaneous variation.

Second, the squared derivative produces
a nonnegative contribution.

Third, the exponential kernel assigns
decreasing weight to older contributions.

The resulting integral is transformed
through a square root and combined with
an initial response value.

The equation therefore defines a
nonnegative memory-dependent functional
of the voltage history.

\subsection{Possible electrical applications}

Potential areas for further investigation include:

\begin{enumerate}

    \item Mathematical analysis of voltage
    variation in time-dependent electrical systems.

    \item Investigation of voltage transient
    histories in circuit models.

    \item Comparison of different voltage
    waveforms through their accumulated
    variation responses.

    \item Mathematical investigation of
    historical voltage dependence.

    \item Development of experimental methods
    for estimating the model parameters.

\end{enumerate}

These applications are prospective.

The functional has not been established
as a direct measurement of physical resistance,
component degradation, or circuit instability.

Such interpretations would require independent
physical derivations and experimental validation.

\section{Limitations and Further Investigation}

The mathematical properties established in
this paper do not independently demonstrate
that the proposed response corresponds to
a new physical electrical quantity.

Several questions remain open.

\subsection{Physical interpretation}

The quantity \(\Omega_s(t)\) has units of ohms,
but dimensional consistency alone does not
establish that it represents electrical resistance.

A physical interpretation would require
a relationship between the functional and
measurable circuit behavior.

\subsection{Parameter identification}

The parameters \(\beta\), \(T\), and \(\lambda\)
must be investigated to determine whether
they can be reliably estimated from experimental
voltage histories.

Their practical identifiability has not
been established.

\subsection{Comparison with existing models}

The exponential memory kernel and squared
derivative are established mathematical
constructions.

Therefore, the proposed combination should
be compared with existing integral functionals,
signal-processing methods, and electrical
system models.

Such comparisons are necessary before
any claim of mathematical or practical novelty
can be justified.

\subsection{Experimental validation}

Future investigations could compare the
predicted response against measurements
obtained from controlled voltage signals.

Such experiments would be necessary to
determine whether the functional provides
information beyond existing electrical
measurements.

\section{Conclusion}

This paper introduced the Electrical Variation
Response Equation, a parameterized mathematical
functional constructed from the historical
variation of a voltage signal.

The proposed equation is

\[
\Omega_s(t)=
\sqrt{
\Omega_0^2+
\frac{\beta}{T}
\int_0^t
e^{-\lambda(t-\tau)}
\left(\frac{dV(\tau)}{d\tau}\right)^2
\,d\tau
}.
\]

The analysis established its dimensional
consistency under appropriate parameter units,
nonnegativity, continuity, and boundedness
under bounded voltage-derivative conditions.

An equivalent first-order differential equation
was derived, demonstrating that the historical
integral can be represented through a local
evolution equation.

The investigation also established parameter
sensitivities, voltage scaling properties,
constant-offset invariance, and expressions
for particular voltage signals.

A reconstruction formula was obtained for
the magnitude of the instantaneous voltage
derivative when the response and model
parameters are known.

The proposed functional provides a mathematical
framework for examining accumulated voltage
variation with exponentially decreasing
historical weighting.

However, its physical significance and
practical usefulness remain open questions.

Further work should investigate its relationship
to existing electrical models, parameter
identifiability, and experimental validation.

\section{Acknowledgments}

The author thanks the teachers, mentors,
and individuals who encourage mathematical
exploration, electrical engineering, and
independent scientific investigation.

\end{document}
